\scanpage{31}
% @ 1
\uline{Prop.} $G$ \deleted{\unsure{groupe alg.}} \tikz[baseline=(s.base)]{\node[inner sep=1pt](s){affine};\node[anchor=south east,inner sep=1pt](a) at ([xshift=2pt,yshift=7pt]s.north west){\inserted{groupe \inserted{(monoïde)}}};\draw ([xshift=-8pt,yshift=-4pt]s.south west) to[out=90,in=190] ([xshift=1pt,yshift=-2pt]a.south west)--([xshift=1pt,yshift=-2pt]a.south east) to[out=0,in=-90] ([xshift=3pt,yshift=3pt]a.south east);} sur le corps $k$. Alors

% @ 2–3
a) $G$ \deleted{\unsure{G}} est limite projective \tikz[baseline=(s.base)]{\node[inner sep=1pt](s){de};\node[anchor=south east,inner sep=1pt](a) at ([xshift=2pt,yshift=7pt]s.north west){\inserted{\unsure{filtr. str|filtr. de|filtrante}}};\draw ([xshift=-8pt,yshift=-4pt]s.south west) to[out=90,in=190] ([xshift=1pt,yshift=-2pt]a.south west)--([xshift=1pt,yshift=-2pt]a.south east) to[out=0,in=-90] ([xshift=3pt,yshift=3pt]a.south east);} \tikz[baseline=(s.base)]{\node[inner sep=1pt](s){(monoïdes)};\node[anchor=south east,inner sep=1pt](a) at ([xshift=2pt,yshift=7pt]s.north west){\inserted{groupes}};\draw ([xshift=-8pt,yshift=-4pt]s.south west) to[out=90,in=190] ([xshift=1pt,yshift=-2pt]a.south west)--([xshift=1pt,yshift=-2pt]a.south east) to[out=0,in=-90] ([xshift=3pt,yshift=3pt]a.south east);} affines de t.f. sur $k$, \deleted{\unsure{et}} \deleted{\unsure{b} avec des morphismes}

% @ 4–5
b) Si $G$ est un monoïde \tikz[baseline=(s.base)]{\node[inner sep=1pt](s){affine};\node[anchor=south east,inner sep=1pt](a) at ([xshift=2pt,yshift=7pt]s.north west){\inserted{(\unsure{grpe|groupe})}};\draw ([xshift=-8pt,yshift=-4pt]s.south west) to[out=90,in=190] ([xshift=1pt,yshift=-2pt]a.south west)--([xshift=1pt,yshift=-2pt]a.south east) to[out=0,in=-90] ([xshift=3pt,yshift=3pt]a.south east);} de t.f. sur $k$, on peut trouver un \deleted{\unsure{i|imm.}} entier $n$ et une immersion fermée
% @ 6
\[
 G\mathrel{\tikz[baseline=-.3ex,x=1em,y=1em]{\draw[line width=.45pt] (0,0)--(1,0) (1,-.12)--(1,.12);}}\tikz[baseline=(s.base)]{\node[inner sep=1pt](s){\(\deleted{\unsure{\mathrm{GL}}}M(n)_{k}^{\unsure{\times|x}}\)};\node[anchor=north,inner sep=1pt](a) at ([xshift=12pt,yshift=-15pt]s.south){\inserted{schéma de matrices $n\times n$}};\draw[->] ([xshift=3pt,yshift=-1pt]s.south) to[out=-100,in=170] ([xshift=-12pt,yshift=2pt]a.north);}.\quad
 (\text{resp. }G\to\mathrm{GL}(n))
\]

% @ 7–9
\uline{Dém} $G\longleftrightarrow A$, \deleted{\unsure{et}} et $A=\underset{i}{\underset{\mathrel{\tikz[baseline=-.3ex,x=1em,y=1em]{\draw[line width=.45pt] (0,0)--(1,0) (1,-.12)--(1,.12);}}}{\mathrm{lim}}}\,A_i$, $A_i$ les sous-algèbres de t. f. de $A$. On s'intéresse aux sous-algèbres $B$ qui sont ``st.bles'' par $\unsure{\Pi|\pi}:A\otimes A$, or il est clair que si
% @ 10–12
$M\subset A$ est st.bl sous $\Pi$, alors aussi la sous-algèbre \tikz[baseline=(s.base)]{\node[inner sep=1pt](s){qu'elle};\node[anchor=south east,inner sep=1pt](a) at ([xshift=2pt,yshift=7pt]s.north west){\inserted{$B$}};\draw ([xshift=-8pt,yshift=-4pt]s.south west) to[out=90,in=190] ([xshift=1pt,yshift=-2pt]a.south west)--([xshift=1pt,yshift=-2pt]a.south east) to[out=0,in=-90] ([xshift=3pt,yshift=3pt]a.south east);} engendre. Donc le corollaire 2 nous est que l'on peut supposer les $A_i$ stbles par $\Pi$. Si $G$ un groupe, d'où un antipodisme
% @ 13–17
$\sigma$ de $A$, alors si $M$ st.bl sous $\sigma$, alors l'algèbre $B$ aussi. Or si $M$ stable par $\Pi$, $\sigma M$ aussi [$M$ sous-module à g. \tikz[baseline=(s.base)]{\node[inner sep=1pt](s){sur};\node[anchor=south east,inner sep=1pt](a) at ([xshift=2pt,yshift=7pt]s.north west){\inserted{(à dr)}};\draw ([xshift=-8pt,yshift=-4pt]s.south west) to[out=90,in=190] ([xshift=1pt,yshift=-2pt]a.south west)--([xshift=1pt,yshift=-2pt]a.south east) to[out=0,in=-90] ([xshift=3pt,yshift=3pt]a.south east);} $\check A$ $\Longleftrightarrow$ $\sigma M$ sous-module à dr (g) sur $\check A$] donc $M+\sigma M$ aussi, donc on peut supposer les $A_i$ de plus st.ble par $\sigma$, donc les $G_i$ sont des groupes.

% @ 18–21
Supposons \unsure{maintnt|maintenant} $A$ de t-f., donc on a un sous $k$-module \tikz[baseline=(s.base)]{\node[inner sep=1pt](s){$M$};\node[anchor=south east,inner sep=1pt](a) at ([xshift=2pt,yshift=7pt]s.north west){\inserted{de tf.}};\draw ([xshift=-8pt,yshift=-4pt]s.south west) to[out=90,in=190] ([xshift=1pt,yshift=-2pt]a.south west)--([xshift=1pt,yshift=-2pt]a.south east) to[out=0,in=-90] ([xshift=3pt,yshift=3pt]a.south east);} de $A$ qui engendre $A$ comme algèbre. \deleted{Considérons} On peut supposer $M$ \unsure{st.bl|stbl} par $\check A$ à \tikz[baseline=(s.base)]{\node[inner sep=1pt](s){\deleted{\unsure{droite}}};\node[anchor=south,inner sep=1pt] at ([yshift=3pt]s.north){\deleted{\unsure{gauche|g.}}};\node[anchor=north,inner sep=1pt] at ([yshift=-3pt]s.south){\inserted{droite}};} donc que $G$ y opère à \deleted{\unsure{gauche}} \inserted{droite}. Donc
% @ 22
\[
 G\longrightarrow\deleted{\unsure{\mathrm{GL}|{\mathrm{End}}}}\,\underline{\mathrm{Aut}}_{k}(M)\quad(\simeq\mathrm{GL}(n)_{k})
\]
% @ 23–24
Av $n=\dim M$. Prouvons que c'est un mono. \deleted{\unsure{On a}} Soit \deleted{$\unsure{g}\ \varepsilon(f)\ \unsure{=0}$} Considérons \tikz[baseline=(s.base)]{\node[inner sep=1pt](s){la};\node[anchor=south east,inner sep=1pt](a) at ([xshift=2pt,yshift=7pt]s.north west){\inserted{pour \unsure{tt|$\forall$} $\varphi\in M$}};\draw ([xshift=-8pt,yshift=-4pt]s.south west) to[out=90,in=190] ([xshift=1pt,yshift=-2pt]a.south west)--([xshift=1pt,yshift=-2pt]a.south east) to[out=0,in=-90] ([xshift=3pt,yshift=3pt]a.south east);} fonction sur $G$
% @ 25
\[
 G\longrightarrow\mathbb E\qquad \deleted{\unsure{g}}\quad
 g\longmapsto\varepsilon(g.\varphi)=(g\varphi)(1)=\varphi(g)
\]
% @ 26–28
\deleted{\unsure{donc} $\unsure{E|k}'$} Donc \deleted{\unsure{g}} $M\unsure{\subset|\in}\mathrm{Im}\,[\text{anneau }\tikz[baseline=(s.base)]{\node[inner sep=1pt](s){de};\node[anchor=south east,inner sep=1pt](a) at ([xshift=-2pt,yshift=7pt]s.north west){\inserted{affine}};\draw ([xshift=-12pt,yshift=-2pt]s.north west) to[out=100,in=0] ([xshift=-10pt,yshift=-2pt]a.south east);\draw ([xshift=-12pt,yshift=-2pt]s.north west)--([xshift=-2pt,yshift=10pt]s.north west);}\,\underline{\mathrm{End}}_k(M)\mathrel{\tikz[baseline=-.3ex,x=1em,y=1em]{\draw[line width=.45pt] (0,0)--(1,0) (1,-.12)--(1,.12);}}A]$ \deleted{\unsure{donc}} \unsure{L'anneau|l'anneau} affine de $\underline{\tikz[baseline=(e.base)]{\node[inner sep=0pt](e){\(\mathrm E\)};\draw ([xshift=1pt,yshift=2pt]e.north west) to[out=20,in=140] ([xshift=2pt,yshift=0pt]e.north east);}\mathrm{nd}}_k(M)$ et a fortiori celui de $\mathrm{GL}(n)$ s'en va \uline{sur} $A$, OK.
% archive sheet number: 30
